\begin{definition}[Extractors and dispersers]
Let $\Gamma$ and $\Omega$ be some finite domains.
Let ${\cal C}$ be a class of random variables taking values in  $\Gamma$.
We say a random variable $P$ taking values in $\Omega$ is
\emph{$\eps$-close to uniform} if for every  $A\subseteq \Omega$,
%
%
\[\left|\Pr (P \in A) - \frac{|A|}{|\Omega|}\right| \leq \eps\,.\]
\begin{itemize}
%
%
%
%
\item Fix any $0\leq \eps <1$. A function $E:\Gamma \to\Omega$
is an \emph{$\eps$-extractor for ${\cal C}$}
if for every $X \in {\cal C}$,
the random variable %
$E(X)$ is $\eps$-close to %
uniform.
%
%
\item A function $D:\Gamma \to\Omega$ is
%
a \emph{disperser for ${\cal C}$} if for every $X \in {\cal C}$,
the random variable %
$D(X)$ takes more than %
one value in $\Omega$ with nonzero probability.
\end{itemize}
\end{definition}
\subsection{Polynomial sources}

%
%
%
In this paper we construct extractors for \emph{polynomial sources},
which are distributions that are sampled by applying low-degree
%
polynomials to uniform inputs as defined next.
%
%
%
Throughout this paper, if $\Omega$ is a finite set, we let
$U_\Omega$ denote the uniform distribution over $\Omega$.
%
%
By the \emph{individual degree} of a multivariate polynomial $f$ we
mean the smallest $d$ such that $f$ has degree $\le d$
in each variable.


\begin{definition}[Polynomial sources]\label{dfn-polysource}
Fix integers $n,k,d$ with $k\leq n$ and a field $\F$.
We define $\M{n}{k}{d}$ to be the set of mappings $f:\F^r\to\F^n$,
where $r$ is an integer counting the number of inputs to the source and
\[f(Z_1,\ldots,Z_r)=(f_1(Z_1,\ldots,Z_r),\ldots,f_n(Z_1,\ldots,Z_r))\] such that
\begin{itemize}
\item for every $i\in [n]$, $f_i$ is a polynomial in $\F[Z_1,\ldots,Z_r]$
of individual degree at most $d$.
\item The range
%
%
%
%
%
%
%
of $f$ is of size at least $q^k$. Formally,
\[|\{f(z_1,\ldots,z_r)\mid (z_1,\ldots,z_r)\in \F^r\}| \geq q^k\,.\]
\end{itemize}
%
A \emph{$(n,k,d)$-polynomial source} is a random variable of the form
$f(U_{\F^r})$ for some and $f\in \M{n}{k}{d}$ with $r$ inputs.
(When the parameters $n,k,d$ are clear from
%
%
the context, we shall omit them, and simply use the term
``polynomial source.'')

\end{definition}

%
%
\begin{definition}[Polynomial-source extractors]\label{dfn-polyext}
%
Let $\Omega$ be some finite set.
A function $E:\F^n\to \Omega$  is a \emph{\polyext{k}{d}{D}{\eps}} if for
every $f\in \M{n}{k}{d}$ of total degree at most $D$ and $r$ inputs,
$E(f(U_{\F^r}))$ is $\eps$-close to uniform
(where $U_{\F^r}$ denotes the uniform distribution
%
over $\F^r$).
%
%
%
%
%
\end{definition}
\begin{remark}
A few words are in order regarding \expref{Definition}{dfn-polysource}.
\begin{itemize}
\item
%
%
%
 The number of inputs used by our source, denoted by $r$ in
 \expref{Definition}{dfn-polysource}, does not affect the
%
%
parameters of our extractors, and hence
%
we omit this parameter from the definition of
polynomial sources and extractors.

\item In the context of extractors what might have seemed
more natural is to require the
%
%
%
%
random variable
$f(U_{\F^r})$ to have
\emph{min-entropy}\footnote{The min-entropy of a
%
%
%
random variable $X$ is
the largest %
real number $\ell$
such that for every
%
value $x$
we have %
%
%
%
$\Pr(X=x)\leq 2^{-{\ell}}$.
This is the standard measure of randomness in the context of extractors,
%
originating from Chor and Goldreich~\cite{CG88}.}
at least $k\cdot \log q$.
%
%
%
Our requirement on the size of the range of $f$ is seemingly weaker,
and suffices for our construction to work.
%
%
%
%
(In particular, our result implies that, for some settings of field order and degree, when $f$ has large range
the random variable $f(U_{\F^r})$ is statistically close to a random variable that has at least a certain min-entropy.)

\item Individual degree plays a larger role than total degree in our results.
%
  In fact, the first stage of our construction---constructing a
%
%
  non-constant polynomial over $\F$---requires a field of order
%
  depending \emph{only} on individual degree. This is why it is
  more convenient to limit individual degree and not total degree
  in the definition of $\M{n}{k}{d}$.

\end{itemize}
\end{remark}
